\documentclass{article}

\usepackage{url,algorithm,amsthm,amsmath,amsfonts,verbatim}
\usepackage{stix}
\newtheorem{lem}{Lemma}
\newtheorem{thm}{Theorem}
\usepackage[toc,page]{appendix}
\newcommand{\Dmax}{D^\mathrm{max}}
\newcommand{\Wmin}{W^\mathrm{min}}
% some traditional definitions that can be blamed on craig barratt
\newcommand{\BEAS}{\begin{eqnarray*}}
\newcommand{\EEAS}{\end{eqnarray*}}
\newcommand{\BEQ}{\begin{equation}}
\newcommand{\EEQ}{\end{equation}}
\newcommand{\BIT}{\begin{itemize}}
\newcommand{\EIT}{\end{itemize}}

% text abbrevs
\newcommand{\eg}{{\it e.g.}}
\newcommand{\ie}{{\it i.e.}}
\newcommand{\ones}{\mathbf 1}

% std math stuff
\newcommand{\reals}{{\mbox{\bf R}}}
\newcommand{\symm}{{\mbox{\bf S}}}

% probability stuff
\newcommand{\Expect}{\mathbf{E}}
\newcommand{\prob}{\mathbf{Prob}}
\newcommand{\var}{\mathop{\bf var}}

% convexity & optimization stuff
\newcommand{\Co}{{\mathop {\bf Co}}}
\newcommand{\co}{{\mathop {\bf Co}}}
\newcommand{\Var}{\mathop{\bf var{}}}
\newcommand{\dist}{\mathop{\bf dist{}}}
\newcommand{\argmin}{\mathop{\rm argmin}}
\newcommand{\arginf}{\mathop{\rm arginf}}

\newcommand{\dom}{\mathop{\bf dom}}
\newcommand{\argmax}{\mathop{\rm argmax}}

\newtheorem{theorem}{Theorem}

% alg description --- not much for now!
\newenvironment{algdesc}%
{\begin{quote}}{\end{quote}}

% figures
\def\figbox#1{\framebox[\hsize]{\hfil\parbox{0.9\hsize}{#1}}}

% captions a la sirev
\makeatletter
\long\def\@makecaption#1#2{
   \vskip 9pt
   \begin{small}
   \setbox\@tempboxa\hbox{{\bf #1:} #2}
   \ifdim \wd\@tempboxa > 5.5in
        \begin{center}
        \begin{minipage}[t]{5.5in}
        \addtolength{\baselineskip}{-0.95pt}
        {\bf #1:} #2 \par
        \addtolength{\baselineskip}{0.95pt}
        \end{minipage}
        \end{center}
   \else
        \hbox to\hsize{\hfil\box\@tempboxa\hfil}
   \fi
   \end{small}\par
}
\makeatother

\usepackage[preprint]{neurips_2021}
\usepackage{fontspec}
\usepackage{xeCJK}
\setCJKmainfont{Songti SC}
\setCJKsansfont{PingFang SC}
\setCJKmonofont{Heiti SC}
\XeTeXlinebreaklocale "zh"
\XeTeXlinebreakskip = 0pt plus 1pt
\usepackage{hyperref}
\usepackage{booktabs}
\usepackage{nicefrac}
\usepackage{microtype}
\usepackage{graphicx}
\usepackage{color}

\graphicspath{{figures/build/}}
\hypersetup{
    pdftitle={Trade Intent Models：跨市场交易与 Agent 策略执行的语义接口},
    pdfauthor={Bill Sun and Alpha.Dev team},
    colorlinks=true,
    urlcolor=blue,
    linkcolor=blue,
    citecolor=blue
}

\renewenvironment{abstract}
{%
  \vskip 0.075in%
  \centerline{\large\bf 摘要}%
  \vspace{0.5ex}%
  \begin{quote}%
}
{
  \par%
  \end{quote}%
  \vskip 1ex%
}

\renewcommand{\figurename}{图}
\renewcommand{\tablename}{表}

\title{Trade Intent Models：\\跨市场交易与 Agent 策略执行的语义接口}
\author{Bill Sun and Alpha.Dev team\\ \url{https://alpha.dev/}}
\date{}

\begin{document}
\maketitle
\small

\begin{abstract}
交易系统仍然被资产类别、交易场所、执行方式与技术栈所割裂。一个经济想法往往需要分别映射到券商 API、中心化交易所、去中心化协议、期权市场、预测市场以及链上交易流程，而这些系统各自拥有不同的语法与操作假设。本文主张，应当把 trade intent 视为交易系统中的第一类语义层。所谓 trade intent model，本质上是一种领域专用语言，它把含糊的人类请求编译成精确、可机读验证、且与具体执行后端解耦的经济动作表示。它在系统中的角色，类似于搜索里的 semantic parsing 或优化中的 CVXPY：用户描述想要达成什么，系统将其编译成规范化表示，再由下游模块翻译成可执行指令。本文的核心论点是，一个设计良好的 trade intent model 会在语义精度、跨 venue 互操作性、策略可组合性与 agent 化自动执行四个维度创造价值，并成为更大一套 action-oriented workflow 的表示层。
\end{abstract}

\section{引言}

问题的核心很简单：人的交易意图是高层的，而执行基础设施却是碎片化、低层化的。用户通常以结果、暴露、对冲与约束来思考问题；但交易 venue 暴露给外部的往往是订单类型、产品标识符、链上交易格式、路由规则和互不兼容的 API。把前者直接映射到后者，天然会导致脆弱的系统。

而且，这个问题正在因交易入口的变化而迅速加剧。对于人类用户而言，越来越多的输入并不是标准化订单，而是 chat、voice，甚至可以称为 vibe trading 的模糊指令，例如“帮我在收盘前拿一点上行暴露”“如果市场反转就帮我对冲”“在多个 venue 中买到这个 thesis 的最佳表达”。这些表达对人类来说自然，但并不是安全的执行格式。它们富含意图，却严重缺乏规格。

对于机器端而言，越来越多的用户正在把通用型或非专业 agent 直接连接到钱包、broker API 与 DeFi protocol 上。这些 agent 可以理解自然语言、调用工具、执行多步金融动作，但它们通常缺少一个稳定的语义层，用来隔离“用户真正想表达的意思”和“不可逆的高权限金融执行动作”。在这种情况下，非结构化 prompt 实际上就变成了直接的 action surface。

近期事件说明了这一点。2026 年 2 月，基于 OpenClaw 的 autonomous crypto agent Lobstar Wilde 因误解了一次小额转账请求，向陌生地址误转了约 5240 万枚 LOBSTAR。2026 年 3 月，中国互联网金融协会则警示称，将 OpenClaw 一类系统用于金融场景，可能暴露敏感凭证，并触发错误转账或非预期投资购买。这些案例共同指向一个系统问题：真正缺失的并不总是预测能力，也未必是智能合约本身，而是位于含糊语言与高权限动作之间的严格中间层。

这种错位具体表现为三类持续性失败：自然语言带有歧义，执行接口彼此割裂，而没有稳定中间表示的自动化系统则难以审计、难以调试。

对应的编译流程可以写成：
\begin{center}
\small
用户想法或自然语言指令 $\rightarrow$ 高层 trade intent $\rightarrow$ \\
执行任务与执行策略 $\rightarrow$ venue-specific orders and transactions.
\end{center}
这个框架把目标从“更好地下单”提升为一个面向 action 的语义接口。

\section{Trade Intent 作为领域专用语言}

Trade intent model 是对期望经济动作的形式化中间表示。最恰当的理解方式，是把它视为一种用于交易语义的领域专用语言。它的作用并不只是记录订单，而是以一种精确、规范、与执行后端解耦且可供机器操作的方式，表达用户真正想要的东西。

最关键的设计原则是：schema 应当表达经济意义，而不是 venue 的表层语法。一个有用的 intent object 需要容纳经济目标、工具与动作类型、组合结构、约束条件、执行偏好，以及从 parsing 到 validation 再到 lowering 的 provenance。也正是因为有了这一层，不同表述的相似请求才能映射到接近的结构，下游模块也才能围绕同一个对象进行算法化推理，而不是一条条处理孤立指令。

例如下面这样的自然语言请求：
\begin{quote}
\small
在 Solana 上买 10 美元的 NVDAx；在 Base 与 Arbitrum 中选择流动性更好的 venue 买 5 美元的 PAX Gold；买一张关于“2026 年 2 月 NVDA 能否突破 200 美元”的 Yes 合约，再买一张关于“今年美联储会降息 2 次”的合约；同时通过 Alpaca via MCP 建立一个 6 月 6 日到期、190/200 行权价的 AAPL bull call spread。
\end{quote}
这个请求在经济意义上是连贯的，但在执行层面却高度异构。Trade intent model 的作用，就是把它视为一个结构化 trade vector，而不是一包彼此无关的 API 调用。

\section{为什么 Trade Intent Model 重要}

Trade intent model 至少在四个维度产生价值。第一，它提高语义精度与可调试性，使系统能够在执行前发现信息缺失、约束冲突或不可行请求。第二，它实现跨 venue 互操作性，使现货、期权、perp、预测市场、broker 流程与链上动作都能通过同一语义接口表示。第三，它提升策略可组合性，使 multi-leg package、路由替代方案、成本与风险权衡、以及 portfolio-level constraints 可以联合评估。第四，它为 agent 化自动执行提供共享对象，使 parser、validator、simulator、planner、monitor 与 agent 能围绕同一个 contract 协同工作，而不必依赖脆弱的 prompt glue。

如果缺少这一中间表示，这些优势就会退化成黑箱自动化；而一旦拥有它，失败也就可以被明确定位为语言理解、intent 构造，或执行翻译中的某一层问题。

\section{作为扩展层的分层执行}

执行本身也必须分层。最好的类比来自机器人系统：high-level planning 与 low-level motor execution 是两层不同的问题。交易系统同样应当把 high-level trade intent 与 low-level execution tasks / policies 区分开来。

但在本文语境下，这一分层执行栈应被理解为对当前仓库工作的扩展，而不是对现有代码已经完整实现这些层的陈述。高层 trade intent 是规划层。它以不保留经济性歧义的方式承载用户目标，同时又提供了进行 cross-venue reasoning、liquidity aggregation 与 policy-aware routing 的正确界面。

例如，``在收盘前买入 10 万股 AAPL''，或者 ``本周买入 100 万美元名义的 NVDA call option''，都不应被直接当作单一订单处理，而应先被表示成包含名义规模、时间范围、impact tolerance、紧迫度与 venue constraints 的 portfolio-level object。只有这样，它们才应被进一步分解成低层 execution task，并分别通过 TWAP、VWAP、broker-native smart routing、链上的 naive slicing、AMM / aggregator execution，或者包含 market-structure / macro-sensitive signal 的更复杂低延迟策略来完成。

这种分层也为 microstructure-level improvement 提供了清晰位置。语义层负责 economic correctness，执行层负责 realized execution quality。用户只在自然语言层指定一次 portfolio objective，而系统可以在不改变原始 intent 的前提下，不断搜索更优的低层实现。当前 GitHub 仓库并未实现这一完整的 planning-and-execution hierarchy；它是在本文中被提出和延展的下一层架构。

\section{系统视角与仓库边界}

这一模块的最终目标，不应被理解为一个“更方便的券商界面”，而更像是一层 action-description layer 与 agent workflow substrate：用户只提供一个 idea，系统便能合成候选策略、完成验证，并支持一键部署到 agent platform。

\begin{figure}[t]
    \centering
    \includegraphics[width=0.84\linewidth]{pipeline-figure.pdf}
    \caption{TIM 作为稳定表示层，连接 idea ingestion、execution-policy selection 与 venue-native deployment。}
    \label{fig:pipeline-zh}
\end{figure}

从系统角度看，pipeline 可以概括为：(1) idea ingestion；(2) semantic parsing into high-level trade intent；(3) strategy synthesis and task decomposition；(4) validation and policy checks；(5) low-level execution policy selection；以及 (6) execution compilation and deployment。

Trade intent schema 正是贯穿这些阶段的稳定接口。配套的 GitHub 仓库 \texttt{galpha-ai/intent-kit}（\href{https://github.com/galpha-ai/intent-kit}{github.com/galpha-ai/intent-kit}）应被理解为这一架构中 contract layer 的当前参考实现：schema definition、template generation、parsing、validation 与 dispatch。也正因为如此，parser、schema、validator、simulator、adapter 与 agent 可以独立演化，但仍通过共享表示保持兼容。

为避免过度表述，需要明确区分“当前仓库已实现的部分”和“本文提出的完整系统”。当前仓库覆盖的是 semantic gateway layer：intent schema、agent-facing template、parsing、validation 与 dispatch；而 high-level portfolio planning、strategy synthesis、execution-policy search 与 low-level algorithmic execution 则属于建立在这一接口之上的未来层。

\section{设计原则与开放问题}

一个可落地的 trade intent language 至少应满足五个原则：语义完整性、规范化结构、关注点分离、可验证性与可扩展性。一个 wrapper 只是翻译语法；真正的 trade intent model 组织的是语义。仍然存在若干重要研究问题：schema 的抽象层级如何选择；lowering 过程中如何保证 semantic preservation；在用户请求不完整时哪些字段可以安全推断；partial fill 与 cross-venue failure recovery 如何处理；以及在更 agentic 的系统中，risk、policy 与 audit layer 应如何约束 agent autonomy。这些都不是边缘问题，而是 deployment-grade trade intent system 的核心设计约束。

\section{结论}

Trade intent model 在人类交易想法与异构执行基础设施之间提供了一个语义接口。它把含糊输入转化为精确、可机读验证、且与具体执行后端解耦的经济动作表示。它的价值远不止更整洁的 order syntax：一个设计良好的 trade intent layer 能实现 cross-venue interoperability、multi-leg strategy composability、可调试性提升，以及 high-level economic intent 与 low-level execution policy 的原则性分离。从这个意义上说，trade intent model 不应被视为某个产品 feature，而应被视为一种基础设施：它是让碎片化市场、异构 API 与 agent platform 能通过同一种 economic-intent language 协同工作的共享语义底座。

\end{document}
